\section{A uniform early prefix}
\label{sec:prefix}

This section supplies all the graph properties used in the continuation.
In addition to degrees and link cycles, we need two upper bounds for
small rooted graphs. The second is useful when all the remaining choices
together have expected count below one.

Fix
\begin{equation}\label{eq:constants}
 \epsilon=\frac1{2000},\qquad r_0=8000,\qquad
 L=100,\qquad B=50,\qquad
 H_0=\max\{r_0+1,3L+2\}=8001.
\end{equation}
Here $r_0$ and $H_0$ bound the link-cycle lengths and template orders
controlled by the prefix; $L$ and $B$ are fixed for the collision-witness
count in Section~\ref{sec:witnesses}.
Write $G_i$ for the graph after $i$ removals, provided that step is reached,
and put
\begin{align}
 p_i&=1-\frac1n-\frac{6i}{n^2},\notag\\
 i_0&=\left\lfloor\frac{n^2}{6}
             \left(1-\frac1n-n^{-1/2+\epsilon}\right)\right\rfloor,
 \qquad p=p_{i_0},\qquad D=np^2,\qquad m=\frac{n^2p}{2}.
 \label{eq:prefix-time}
\end{align}
These quantities are deterministic. In particular,
\begin{equation}\label{eq:scales}
 p=n^{-1/2+\epsilon}+O(n^{-2}),\qquad
 D\sim n^{2\epsilon},\qquad e(G_{i_0})=m
\end{equation}
whenever $i_0$ is reached.

\begin{definition}[Rooted templates]\label{def:template}
A template $(H,I)$ consists of a graph $H$ and an independent set
$I\subseteq V(H)$ of distinguished vertices. For an injection
$\psi:I\to[n]$, let $X_{H,I,\psi}(G)$ count its edge-preserving injective
extensions to $V(H)$ in $G$. Its scaling at density $p$ is
\[
 S(H,I)=n^{v(H)-|I|}p^{e(H)}.
\]
A subtemplate with the same distinguished set is any subgraph
$H'\subseteq H$ whose vertex set contains $I$, distinguished by $I$.
When $v(H)>|I|$, the rooted density is $e(H)/(v(H)-|I|)$; it is zero
when there are no free vertices. The template is \emph{balanced} if
every same-root subtemplate has rooted density at most that of $(H,I)$.
\end{definition}

The independence of $I$ is a property of the template. Its prescribed
images may have additional edges in $G$.

\begin{proposition}[Good prefix]\label{prop:prefix}
There are constants $c>0$ and $C_0<\infty$ and an event $\mathcal G_n$
measurable from the process through step $i_0$ such that
\[
 \Prb(\mathcal G_n^c)\le
       \exp\bigl(- (\log n)^{4/3}\bigr)=n^{-\omega(1)}
\]
for all sufficiently large $n$. On this event, $i_0$ is reached and
$G=G_{i_0}$ has $m$ edges and the following simultaneous properties.
\begin{enumerate}[label=\textup{(\roman*)}]
\item\label{it:degrees}
Every vertex has degree $(1\pm n^{-c})np$, and every two distinct
vertices have $(1\pm n^{-c})D$ common neighbors.
\item\label{it:cycles}
For every vertex $u$ and $3\le j\le r_0$, the number of embeddings of a
fixed labeled cycle $C_j$ into the link $G[N_G(u)]$ is
$(1\pm n^{-c})D^j$.
\item\label{it:templates}
For every template $(H,I)$ with $v(H)\le H_0$ and every injection
$\psi:I\to[n]$, the following two bounds hold.
If $S(H',I)\ge1$ for every same-root subtemplate $H'$, then
\begin{equation}\label{eq:template-large}
 X_{H,I,\psi}(G)\le(1+\log n)^{C_0}S(H,I).
\end{equation}
If $n^{v(H)-|J|}p^{e(H)-e(H[J])}\le1$ for every
$I\subseteq J\subseteq V(H)$, then
\begin{equation}\label{eq:template-small}
 X_{H,I,\psi}(G)\le(1+\log n)^{C_0}.
\end{equation}
\end{enumerate}
All constants and bounds are uniform over the graph and the prescribed
root images.
\end{proposition}

\begin{proof}
We give the specialization of Joos--K\"uhn's stopping estimates, including
the parameters and rooted conventions. Their Sections~5--7 (pp.~9--11)
set up the trajectories and stopping times; Lemma~10.1 (p.~50) controls
those times, and Lemma~7.15 (pp.~16--17) supplies the two general extension bounds
\cite{joosKuehn2025}.

\paragraph{Public parameters and initial hypotheses.}
Denote their small parameter by $\eta$, keeping it distinct from
$\epsilon$ in \eqref{eq:constants}. Choose $\eta$ small enough for their
hierarchy and so that
\[
 \eta<\epsilon/2,\qquad \eta^{-4}\ge H_0.
\]
Choose their $\delta>0$ sufficiently small in terms of $\eta$, with
$\delta<1/2$. These choices remain fixed as $n$ grows.
The removed graph is a triangle, with uniformity parameter $k=2$ and
2-density $(3-1)/(3-2)=2$; it is strictly 2-balanced.
The initial density in their normalization is
$\theta=2e(K_n)/n^2=1-1/n$. Their density trajectory is therefore exactly
$p_i$ in \eqref{eq:prefix-time}. Their horizon has density
$n^{-1/2+\eta}$, strictly below $p_{i_0}$ for large $n$.

Their Section~5 requires an $(\eta^4,\delta,2)$-pseudorandom initial
graph. For completeness, $K_n$ satisfies all its defining conditions
(P1)--(P4) on p.~3. The public convention permits edges among
distinguished vertices and ignores them in both counts and scalings.
For a template with $r$ distinguished vertices, $a$ other vertices, and
$b$ edges not internal to the distinguished set, every fixed root
injection has exactly $(n-r)_a$ extensions in $K_n$; here
$(x)_a=x(x-1)\cdots(x-a+1)$ and $(x)_0=1$.
For all templates with at most $\eta^{-4}$ vertices,
\[
 \frac{(n-r)_a}{n^a\theta^b}=1+O_\eta(n^{-1}).
\]
Their initial error parameter is
$\zeta_{\mathrm{init}}=n^\delta/(\sqrt n\,\theta)$.
The displayed relative error is smaller than both
$\zeta_{\mathrm{init}}$ and $\zeta_{\mathrm{init}}^\delta$, which verifies
the relative-error requirements (P1) and (P2).
If $a\ge1$, the scaling is asymptotic to $n^a$, while the upper scaling
threshold in (P3) is at most $n^{1/2}$; thus neither small-scaling condition
(P3) or (P4) applies. If $a=0$, count and scaling are both exactly one,
and the remaining conditions hold directly.

\paragraph{Balanced counts and the stopping threshold.}
The dynamic error parameter is
\[
 \zeta_i=\frac{n^{\eta^2}}{\sqrt n\,p_i}
         =n^{\eta^2}x_i^{-1/2},\qquad x_i=np_i^2.
\]
The balanced-template controls in Section~7 give relative error
$\zeta_i^\delta$ until the relevant scaling reaches
$\zeta_i^{-\sqrt\delta}$, simultaneously over the templates and root
injections in their families. For a balanced rooted template with
$a\ge1$ free vertices and density at most two, its scaling is at least
$x_i^a$. Throughout $i\le i_0$, we have $x_i\ge n^{2\epsilon}$ and
\[
 \frac{x_i^a}{\zeta_i^{-\sqrt\delta}}
  =n^{\eta^2\sqrt\delta}x_i^{a-\sqrt\delta/2}
  \longrightarrow\infty.
\]
Thus each template used for parts~\ref{it:degrees} and~\ref{it:cycles}
belongs to the tracked family initially and never reaches its threshold
before $i_0$.

A rooted edge counts degrees. A path of length two with its endpoints
distinguished counts common neighbors, even when the prescribed endpoint
images are adjacent. For a link cycle of length $j$, use the wheel whose
center is distinguished and whose rim is the labeled $C_j$. This template
has $j$ free vertices and $2j$ edges. Any subset of $s$ free vertices has
at most $s$ spokes and $s$ rim edges, so the wheel is balanced with density
two. Its extensions are exactly the labeled link-cycle embeddings and
have scaling $D^j$. The largest wheel has $r_0+1\le\eta^{-4}$ vertices.
The degree and codegree templates are likewise balanced, with scalings
$np$ and $D$ respectively.

At $i_0$,
\[
 \zeta_{i_0}\le n^{-(\epsilon-\eta^2)}.
\]
Choose any fixed $c$ with
$0<c<\delta(\epsilon-\eta^2)$. The public relative errors then imply
parts~\ref{it:degrees} and~\ref{it:cycles}, with slack for fixed constants.

\paragraph{General rooted extensions.}
On the same stopping controls, Lemma~7.15 applies to every template up
to the public size cap, with no balance assumption on the target
template. Its first case states that, if all same-root subtemplate
scalings are at least one, the extension count is at most its scaling
times $(1+\log n)^{\alpha_{H,I}}$. Its second case gives the bound
$(1+\log n)^{\alpha_{H,I}}$ when every completion scaling after enlarging
the distinguished set is at most one. Here
\[
 \alpha_{H,I}=2^{v(H)-|I|+1}-2.
\]
In our independent-root convention these hypotheses are exactly those
in \eqref{eq:template-large} and \eqref{eq:template-small}. For the latter,
the public convention ignores the edges internal to the enlarged root
set $J$, giving precisely the exponent $e(H)-e(H[J])$.
We may take $C_0=2^{H_0+1}-2$. Its size is immaterial: it is fixed
independently of $n$.

\paragraph{Probability and conditioning.}
Lemma~10.1 bounds the probability of failure of the required stopping
controls before the public horizon by
$\exp(- (\log n)^{4/3})$. In particular the triangle count is controlled
through $i_0$ around the positive trajectory $n^3p_i^3/6$, with relative
error tending to zero. Termination at or before $i_0$ would contradict
that control, so $i_0$ is reached on this event.
Define $\mathcal G_n$ to be the event of reaching $i_0$ and satisfying
the properties in the proposition at that time. It is determined by
the prefix alone. The later public success event is a subset of
$\mathcal G_n$ and proves its probability estimate. We do not condition
the continuation on any future stopping control.
\end{proof}

From now through Section~\ref{sec:witnesses}, fix any prefix in
$\mathcal G_n$ and write $G=G_{i_0}$. All estimates will be uniform over
this choice. Let $\cE=E(G)$, and let $\cT$ be the set of triangles of
$G$, each regarded as a three-element subset of $\cE$. Thus $|\cE|=m$,
and distinct members of $\cT$ have at most one common edge. For
$e\in\cE$, let
\[
 d_e=|\{T\in\cT:e\in T\}|=(1\pm n^{-c})D.
\]
This last equality is the common-neighbor estimate for the endpoints
of $e$. The task is now to analyze the removal process started from
this fixed graph.
